%! AP Statistics — Unit 4, Lesson 4.3 — Group Activity
\documentclass[10pt]{article}
\usepackage{apstats-article}
\usepackage{apstats-boxes}

\begin{document}

\pageheader{Unit 4, Lesson 4.3}{Group Activity: Cards, Spinners, and the Language of Probability}
\namepartnerperiod

\vspace{0.1in}

\begin{scenariobox}[Context]{sky}
\small
A standard deck of 52 playing cards has 4 suits (hearts, diamonds, clubs, spades),
13 cards per suit (Ace, 2--10, Jack, Queen, King). Hearts and diamonds are red;
clubs and spades are black. Face cards = Jack, Queen, King (3 per suit, 12 total).
Work through the tier your teacher assigns.
\end{scenariobox}

\vspace{0.1in}

% ── Tier R ────────────────────────────────────────────────────────────────────
\begin{tcolorbox}[colback=goldbg, colframe=goldacc, boxrule=1pt, arc=2mm,
    left=4mm, right=4mm, top=3mm, bottom=3mm,
    title={\bfseries Tier R \enspace Remediate --- Building Sample Spaces}]
\small

One card is drawn at random from a standard 52-card deck.

\begin{enumerate}[leftmargin=1.5em, itemsep=10pt]
  \item How many total outcomes are in the sample space? \blank{2cm}

  \item Calculate each probability. Show work as a fraction.
        \vspace{4pt}
        \renewcommand{\arraystretch}{1.8}
        \begin{tabularx}{\linewidth}{@{} l X @{}}
        Event & $P(E) =$ \\
        \hline
        Drawing a heart & \blank{4cm} \\
        Drawing a red card & \blank{4cm} \\
        Drawing a face card & \blank{4cm} \\
        Drawing the Ace of Spades & \blank{4cm} \\
        \end{tabularx}

  \item Apply the complement rule:
        \begin{enumerate}[label=(\alph*), itemsep=6pt]
          \item $P(\text{not a heart}) = \blank{3cm}$
          \item $P(\text{not a face card}) = \blank{3cm}$
        \end{enumerate}

  \item Is it possible to draw a card with probability 0? Probability 1? Explain.\\
        \writeline\writeline
\end{enumerate}

\end{tcolorbox}

\vspace{0.1in}

% ── Tier A ────────────────────────────────────────────────────────────────────
\begin{tcolorbox}[colback=greenbg, colframe=greenacc, boxrule=1pt, arc=2mm,
    left=4mm, right=4mm, top=3mm, bottom=3mm,
    title={\bfseries Tier A \enspace Approaching Proficiency --- Probability and Interpretation}]
\small

One card is drawn at random from a standard 52-card deck.

\begin{enumerate}[leftmargin=1.5em, itemsep=10pt]
  \item Calculate and interpret each probability in the context of the card draw.
        \vspace{4pt}
        \renewcommand{\arraystretch}{1.8}
        \begin{tabularx}{\linewidth}{@{} l p{2.5cm} X @{}}
        Event & $P(E)$ & Long-run interpretation \\
        \hline
        Drawing a club & \blank{2.2cm} & \writeline \\
        Drawing a King & \blank{2.2cm} & \writeline \\
        Drawing a red face card & \blank{2.2cm} & \writeline \\
        \end{tabularx}

  \item A friend says: ``I drew a heart, so the probability of drawing a heart is 1.''
        Explain the error in this reasoning.\\
        \writeline\writeline

  \item Use the complement rule to find:
        \begin{enumerate}[label=(\alph*), itemsep=6pt]
          \item $P(\text{not a King}) = \blank{3cm}$
          \item $P(\text{not a red face card}) = \blank{3cm}$
        \end{enumerate}

  \item Use a \textbf{tree diagram} to list the sample space for drawing the suit and
        color of one card. (Hint: there are 2 colors and 4 suits; which suits go with
        which color?) Label each branch and compute $P(\text{red})$ from the tree.
        \vspace{50pt}
\end{enumerate}

\end{tcolorbox}

\vspace{0.1in}

% ── Tier E ────────────────────────────────────────────────────────────────────
\begin{tcolorbox}[colback=sky, colframe=navy, boxrule=1pt, arc=2mm,
    left=4mm, right=4mm, top=3mm, bottom=3mm,
    title={\bfseries Tier E \enspace Extension --- Multi-Stage Sample Spaces}]
\small

A spinner has 10 equally likely sections: sections 1--5 are red, sections 6--8 are
blue, and sections 9--10 are green. The spinner is spun \textbf{twice}.

\begin{enumerate}[leftmargin=1.5em, itemsep=10pt]
  \item How many outcomes are in the sample space for two spins? \blank{3cm}
  \item Use a tree diagram or organized list to count the number of outcomes in which
        \textbf{both spins land on red}.
        \vspace{40pt}
  \item Calculate $P(\text{both red}) = \blank{3cm}$
  \item Calculate $P(\text{at least one blue spin}) $ using the complement rule:\\
        Hint: first find $P(\text{no blue on either spin})$.
        \vspace{40pt}
  \item Interpret $P(\text{at least one blue spin})$ in one sentence using the phrase
        ``in the long run.''\\
        \writeline\writeline
  \item (Challenge) If you know that the first spin was red, does the probability that
        the second spin is also red change? Why or why not?\\
        \writeline\writeline
\end{enumerate}

\end{tcolorbox}

\end{document}
